\section{Introduction}

Multimodal semantic segmentation improves scene understanding by combining complementary sensing cues such as RGB, depth, thermal, and other auxiliary modalities. This direction has become increasingly important for robust perception under adverse illumination, occlusion, and environmental uncertainty~\cite{zhang2021multimodalsurvey,hazirbas2016fusenet,zhang2023cmx}. For drone platforms, however, multimodal segmentation is particularly challenging. First, the bird's-eye viewpoint dramatically reduces object scale, making small categories easy to lose during repeated downsampling and decoding. Second, changes in altitude, orientation, and lighting conditions make the relative usefulness of RGB and thermal modalities highly dynamic. As a result, fixed fusion strategies often fail to exploit cross-modal complementarity in a scene-dependent manner. Existing UAV-oriented datasets and models, such as UAVid, Agriculture-Vision, SynDrone, UAVformer, and 2DSegFormer, also indicate that aerial scene parsing differs substantially from ground-view segmentation in scale distribution, long-range context, and structural recovery~\cite{lyu2020uavid,chiu2020agriculturevision,rizzoli2023syndrone,zhang2023uavformer,li2022twodsegformer}.

Recent foundation models have substantially strengthened visual representation learning. Large-scale pre-training with ViT, MAE, and ConvNeXt has established powerful generic backbones~\cite{dosovitskiy2021vit,he2022mae,liu2022convnext}. Building on this trend, SAM and SAM~2 provide strong generic segmentation priors, while ViT-Adapter, SAM-Adapter, SAM2-Adapter, and MM SAM-Adapter demonstrate that frozen foundation backbones with lightweight adaptation modules form an effective route to dense prediction and multimodal transfer~\cite{kirillov2023segmentanything,ravi2025sam2,chen2023vitadapter,chen2023samadapter,chen2024sam2adapter,curti2025mmsamadapter}. At the same time, LoRA-style parameter-efficient tuning strategies and their vision variants further support the design choice of keeping most pre-trained parameters frozen while adapting only a small trainable subset~\cite{hu2022lora,zhong2024convlora,xin2024peftsurvey}. Nevertheless, for drone-view RGB-T segmentation on FMB, two limitations remain clear. Existing methods usually do not explicitly model modality-quality variation at the input-feature level, and their decoders are rarely designed around the severe small-object degradation caused by aerial viewpoints.

To address these limitations, we propose \method{}, a foundation-model-driven RGB-T segmentation framework tailored to drone scenes. The framework uses a frozen SAM2 Hiera-L RGB backbone with bottleneck adapters and a fully trainable ConvNeXt-Tiny thermal encoder in an asymmetric dual-stream design. On top of these encoders, we introduce Condition-Aware Cross-Modal Adaptive Fusion (CACAF), which estimates modality quality and performs adaptive cross-modal enhancement at multiple scales. We further propose the Hierarchical Guided Small-Object-Aware Decoder (HGSOAD), which combines hierarchical decoding, scale-aware gating, and auxiliary supervision to improve the recovery of fine-grained structures and distant small objects. The overall design is centered on FMB, a full-time RGB-T benchmark that is especially suitable for evaluating robustness to modality-quality fluctuations in challenging aerial conditions~\cite{liu2023fmb,curti2025mmsamadapter}.

Our main contributions are summarized as follows:
\begin{enumerate}
    \item We propose CACAF, a condition-aware fusion module that explicitly estimates modality quality and performs adaptive cross-modal enhancement instead of relying on static multimodal fusion.
    \item We propose HGSOAD with SAGU, a lightweight decoder enhancement tailored to the small-object recovery problem in drone-view segmentation.
    \item On the FMB test set, the proposed method achieves \textbf{68.62 mIoU}, improving over MM SAM-Adapter by \textbf{+2.52}.
    \item We provide unified module and loss-function ablations, clarifying the individual roles of CACAF, SAGU, and the CE + Dice training objective.
\end{enumerate}
